\documentclass[10pt,a4paper]{article}
\usepackage[margin=0.8in]{geometry}
\usepackage{amsmath,amssymb,hyperref}
\hypersetup{colorlinks=true,urlcolor=blue,linkcolor=blue}
\setlength{\parindent}{0pt}\setlength{\parskip}{6pt}
\begin{document}
\begin{center}
{\Large The printed inverse-length frequency rule is impossible}\\[4pt]
TLMC 00000000079 \quad lizaixi01 \quad 24 October 2026
\end{center}

\textbf{Exact scope.}
Both versions of the source say that every digit block occurs with frequency equal
to the inverse of its length. This submission disproves that explicit necessary
clause. It does \emph{not} assert that the usual, correctly normalized normality
statement is false. In ordinary base-$b$ normality the prescribed frequency for
a length-$\ell$ word is $b^{-\ell}$, not $\ell^{-1}$; see \cite{ajm},
Definition 1.1 on printed p.~1024. The two requirements already disagree for one digit.

\textbf{The actual irrational and prime sequence.}
Take $\alpha=\sqrt2$ and let $p_n$ be the increasing enumeration of all primes,
starting with $p_0=2$. Define
\[
 x_n=\{\alpha p_n\}=\alpha p_n-\lfloor\alpha p_n\rfloor .
\]
The formal proof uses the native real square root, the native irrationality
theorem for $\sqrt2$, the genuine function \texttt{Nat.nth Nat.Prime}, and
\texttt{Int.fract}. It proves that each $p_n$ is prime, that the enumeration is
strictly increasing, and that $0\le x_n<1$ for every $n$.

\textbf{Decimal digits and point-sequence cylinders.}
Choose base $10$. For $x\in[0,1)$ set
\[
 d_j(x)=\left\lfloor 10\{10^j x\}\right\rfloor\in\{0,\ldots,9\},\qquad j\ge0.
\]
This fixes the terminating expansion at decimal boundary points. For a word
$w=(w_0,\ldots,w_{\ell-1})$, let
\[
 C_N(w)=\#\{0\le n<N:\ d_j(x_n)=w_j\text{ for all }0\le j<\ell\},
 \qquad f_N(w)=\frac{C_N(w)}{N}\quad(N>0).
\]
Thus we count the points whose decimal expansions begin with the word; no limiting
frequency is assigned by definition. In particular the proof establishes the actual
first-digit cylinder identity
\[
 d_0(x)=a \quad\Longleftrightarrow\quad
 \frac a{10}\le x<\frac{a+1}{10}\qquad(0\le a<10).
\]
The value at $N=0$ is immaterial to the limit at infinity; the formal division
convention gives zero, and the counting bound below also holds there.

\textbf{A finite counting obstruction.}
Write $(0)$ and $(1)$ for the two distinct one-digit words. A point cannot begin
with both digits, so the two actual finite sets of indices are disjoint subsets of
$\{0,\ldots,N-1\}$. Therefore, for every $N>0$,
\[
 C_N((0))+C_N((1))\le N,
 \qquad f_N((0))+f_N((1))\le1. \tag{1}
\]
This is a theorem about the defined index sets and their cardinalities. It does not
use any assumption about equidistribution of primes.

\newpage
\textbf{The limiting contradiction.}
The printed inverse-length rule requires, for each of the two length-one words,
\[
 \lim_{N\to\infty} f_N((0))=1,
 \qquad \lim_{N\to\infty} f_N((1))=1.
\]
The sum would tend to $2$, whereas (1) forces its limit to be at most $1$.
This contradiction proves that the actual sequence $x_n=\{\sqrt2 p_n\}$ fails
the printed rule. In fact the proof works for every sequence of real points.
The concrete irrational and prime sequence are included so that the counterexample
satisfies the source's specified objects and its universal irrational quantifier.

\textbf{The digit-stream interpretation gives the same obstruction.}
One can instead form the actual digit stream $a_n=d_0(x_n)$ and count a word at
successive positions:
\[
 S_N(w)=\#\{0\le n<N:\ a_{n+j}=w_j\text{ for all }0\le j<\ell\}.
\]
The formal development also defines these sliding occurrences. For a one-digit
word they coincide exactly with $C_N(w)$, so the same contradiction rules out
the inverse-length requirement for this interpretation. The proof therefore
does not depend on choosing between prefix-cylinder sampling and the sliding
first-digit stream. It makes no claim about other unprinted constructions
of a digit string from the point sequence.

\textbf{Formal coverage and verification.}
\texttt{Main.lean} defines the decimal digits using real floor and fractional part,
the finite occurrence sets using \texttt{Finset.range} and filtering, their actual
cardinalities, and the normalized real frequencies. It proves the first-digit
interval identity, disjointness, the finite inequality, and the limit contradiction.
Its final existential theorem supplies the genuine irrational $\sqrt2$, its native
prime-point sequence in $[0,1)$, and failure of both printed inverse-length rules.
No desired frequency, cardinality, or failure is assigned as a definition.

The companion project pins Lean 4.33.0 and Mathlib
\texttt{db584cd6d46c92f209a44c0f1c829460d327499d}.
It contains no admitted proof, extra axiom or \texttt{native\_decide}.
The submission evidence records the actual build, all-namespace ordinary
\texttt{\#print axioms} audit and official fresh-kernel replay, with source hashes
and process records. Production timings are not controlled efficiency evidence.

\textbf{What the result does not say.}
The correct decimal normalization for a one-digit word is $1/10$.
Then the two relevant frequencies would sum to $1/5$, which is entirely
consistent with (1). This submission neither proves nor refutes the conventional
normality or equidistribution assertion for prime multiples. It only refutes
the literal inverse-of-length clause printed in this problem.

\begin{thebibliography}{9}
\bibitem{ajm} D. Airey, S. Jackson and B. Mance,
\emph{Descriptive complexity in Cantor series},
The Journal of Symbolic Logic \textbf{87} (2022), no.~3, 1023--1045.
\href{https://doi.org/10.1017/jsl.2021.77}{doi:10.1017/jsl.2021.77}.
Definition 1.1 on printed p.~1024 is used solely to distinguish ordinary
normality's $b^{-\ell}$ normalization from the problem's printed $\ell^{-1}$.
\end{thebibliography}
\end{document}
